% !TeX root = ../algebralineare.tex

\section{Strutture algebriche} \label{sec:algstructures}

\subsection{Gruppi} \label{subsec:groups}
    Un gruppo è una struttura algebrica determinata da un insieme $S$ e da un'operazione binaria interna e chiusa $+ : S \times S \rightarrow S$, la quale deve soddisfare una serie
    di assiomi:
    \begin{enumerate}[label=(\arabic*)]
        \item Associatività: $\forall x, y, z \in S, \; x + (y + z) = (x + y) + z$
        \item Esistenza dell'elemento neutro: $\exists 0_G \in S \; | \; \forall x \in S, \; x + 0_G = x$
        \item Esistenza dell'elemento inverso: $\forall x \in S, \exists (-x) \in S \; | \; x + (-x) = 0_G$
    \end{enumerate}
    Qualora un gruppo soddisfi anche l'assioma di commutatività, ovvero $\forall x, y \in S, \; x + y = y + x$, esso viene detto ``gruppo abeliano''.

    \vspace{10pt}
    Alcuni esempi di gruppi sono:
    \begin{itemize}
        \item $(\mathbb{Z}, +)$
        \item $(Q_8, +)$
        \item $(S_n, +)$
    \end{itemize}



\subsection{Anelli} \label{subsec:rings}
    Un anello è una struttura algebrica determinata da un insieme $S$ e da due operazioni binarie interne e chiuse, $+: S \times S \rightarrow S$ (addizione) e $\cdot : S \times S \rightarrow S$ (moltiplicazione), le quali devono soddisfare una serie di assiomi:
    \begin{enumerate}[label=(\arabic*)]
        \item $(S, +)$ deve essere un gruppo abeliano
        \item Associatività della moltiplicazione: $\forall x, y, z \in S, \; x \cdot (y \cdot z) = (x \cdot y) \cdot z$
        \item Esistenza dell'elemento neutro della moltiplicazione: $\exists 1_R \in S \; | \; \forall x \in S, \; x \cdot 1_R = x$
        \item Distributività destra della moltiplicazione rispetto alla somma: $\forall x, y, z \in S, \; x \cdot (y + z) = x \cdot y + x \cdot z$
        \item Distributività sinistra della moltiplicazione rispetto alla somma: $\forall x, y, z \in S, \; (y + z) \cdot x = y \cdot x + z \cdot x$
    \end{enumerate}
    E' necessario sottolineare la differenza fra distributività destra e sinistra in quanto la $\textit{moltiplicazione}$ in un anello non è detto che sia commutativa.

    \vspace{10pt}
    Alcuni esempi di anelli sono:
    \begin{itemize}
        \item $(\mathbb{R}^{n, n}, +, \cdot)$ dove $\cdot$ è la moltiplicazione fra matrici
        \item $(\mathbb{Z}/n\mathbb{Z}, +, \cdot)$
        \item $(\mathbb{R}[x], +, \cdot)$
    \end{itemize}



\subsection{Campi} \label{subsec:fields}
    Un campo è una struttura algebrica determinata da un insieme $S$ e da due operazioni binarie interne e chiuse, $+ : S \times S \rightarrow S$ (addizione)  e $\cdot : S \times S \rightarrow S$ (moltiplicazione), le quali devono soddisfare una serie di assiomi:
    \begin{enumerate}[label=(\arabic*)]
        \item $(S, +)$ deve essere un gruppo abeliano
        \item $(S^*, \cdot)$ deve essere un gruppo abeliano, dove $S^*$ non è altro che $S$ senza gli elementi per cui l'inverso moltiplicativo ($a \in S \implies \exists a^{-1} \in S \; | \; aa^{-1} = \textnormal{id}_{F}(a)$) non è definito ($S / \{0\}$)
        \item Distributività completa della moltiplicazione rispetto alla somma: $\forall x, y, z \in S, \; x \cdot (y + z) = x \cdot y + x \cdot z = (y + z) \cdot x$
    \end{enumerate}

    \vspace{10pt}
    Alcuni esempi di campi sono:
    \begin{itemize}
        \item $(\mathbb{Q}, +, \cdot)$
        \item $(\mathbb{R}, +, \cdot)$
        \item $(\mathbb{C}, +, \cdot)$, il quale è il più piccolo campo algebricamente chiuso.
    \end{itemize}



\subsection{Omomorfismi} \label{subsec:homomorphisms}
    Un omomorfismo è un'applicazione che, associando due strutture algebriche, ne mantiene le proprietà.
    Per esempio, siano $(G, +)$ e $(H, *)$ due gruppi,
    allora $f: G \rightarrow H$ è un omomorfismo se e solo se
    \[
        \forall x, y \in G, \; f(x + y) = f(x) * f(y)
    \]
    Un esempio particolare di omomorfismo sono le applicazioni lineari fra spazi vettoriali $f: V \rightarrow W$ in quanto $f(\mathbf{v} + \mathbf{u}) = f(\mathbf{v}) + f(\mathbf{u})$.
